\subsection{Constructions Derived from the Toolkit}

The following objects are not added as new primitive permissions. Each is
constructed from lines, circles, and their intersections.

\subsubsection{The Perpendicular Bisector and the Midpoint}

\begin{derivationbox}[Construction]
Let $A$ and $B$ be distinct points. Draw two circles with the same radius,
one centred at $A$ and one centred at $B$, choosing the radius larger than half
of $|AB|$. Let their intersection points be $P$ and $Q$.

Because $P$ lies on both circles,
\[
|PA|=|PB|,
\]
and similarly,
\[
|QA|=|QB|.
\]
Thus both $P$ and $Q$ are equidistant from $A$ and $B$.
\end{derivationbox}

\begin{center}
\begin{tikzpicture}[scale=0.85]
    \coordinate (A) at (-2.2,0);
    \coordinate (B) at (2.2,0);
    \coordinate (P) at (0,2.4);
    \coordinate (Q) at (0,-2.4);
    \draw[helper] (A) circle (3.25);
    \draw[helper] (B) circle (3.25);
    \draw[subset] (A) -- (B);
    \draw[very thick] (0,-3.0) -- (0,3.0);
    \fill[geometricSubsetColor] (A) circle (2pt) node[below] {$A$};
    \fill[geometricSubsetColor] (B) circle (2pt) node[below] {$B$};
    \fill[geometricSubsetColor] (P) circle (2pt) node[above] {$P$};
    \fill[geometricSubsetColor] (Q) circle (2pt) node[below] {$Q$};
\end{tikzpicture}
\end{center}

\begin{theorembox}
The line $PQ$ is the perpendicular bisector of $AB$. Its intersection with
$AB$ is the midpoint of $AB$. Thus both midpoint and perpendicularity are
constructed from equal radii and intersections.
\end{theorembox}

\subsubsection{A Perpendicular Through a Point on a Line}

\begin{derivationbox}[Construction]
Let $P$ lie on a line $\ell$.

\textbf{Step 1.} Draw a circle centred at $P$. It meets $\ell$ at two points
$A$ and $B$. Since $|PA|=|PB|$, the point $P$ is the midpoint of $AB$.

\textbf{Step 2.} Apply the perpendicular-bisector construction to $AB$. The
resulting perpendicular bisector passes through its midpoint $P$.
\end{derivationbox}

\begin{center}
\begin{tikzpicture}[scale=0.9]
    \coordinate (A) at (-2.2,0);
    \coordinate (P) at (0,0);
    \coordinate (B) at (2.2,0);
    \coordinate (Q) at (0,2.4);
    \coordinate (R) at (0,-2.4);

    \draw[subset] (-3.5,0) -- (3.5,0) node[right] {$\ell$};
    \draw[dashed,line width=0.9pt,draw=blue!65] (P) circle (2.2);
    \draw[dashed,line width=0.9pt,draw=orange!85!black] (A) circle (3.25);
    \draw[dashed,line width=0.9pt,draw=orange!85!black] (B) circle (3.25);
    \draw[very thick] (Q) -- (R) node[below] {$m=QR$};
    \draw (0.28,0) -- (0.28,0.28) -- (0,0.28);

    \fill[geometricSubsetColor] (A) circle (2pt) node[below] {$A$};
    \fill[geometricSubsetColor] (P) circle (2pt) node[below right] {$P$};
    \fill[geometricSubsetColor] (B) circle (2pt) node[below] {$B$};
    \fill[geometricSubsetColor] (Q) circle (2pt) node[above] {$Q$};
    \fill[geometricSubsetColor] (R) circle (2pt) node[below] {$R$};

    \node[blue!65,align=left] at (-3.0,2.75)
        {\small Step 1: circle centred at $P$};
    \node[orange!85!black,align=left] at (3.0,-2.75)
        {\small Step 2: circles centred at $A$ and $B$};
\end{tikzpicture}
\end{center}

\begin{theorembox}
The constructed line $m$ passes through $P$ and is perpendicular to $\ell$.
\end{theorembox}

\subsubsection{A Perpendicular Through a Point Outside a Line}

\begin{derivationbox}[Construction]
Let $P$ lie outside a line $\ell$.

\textbf{Step 1.} Draw a circle centred at $P$ with radius large enough to meet
$\ell$ at two points $A$ and $B$. Then $|PA|=|PB|$, so $P$ lies on the
perpendicular bisector of $AB$.

\textbf{Step 2.} Draw equal-radius circles centred at $A$ and $B$, chosen so that
they meet again at a second point $Q$. The points $P$ and $Q$ are both
equidistant from $A$ and $B$.
\end{derivationbox}

\begin{center}
\begin{tikzpicture}[scale=0.9]
    \coordinate (P) at (0,2.0);
    \coordinate (Q) at (0,-2.0);
    \coordinate (A) at (-2.1,0);
    \coordinate (B) at (2.1,0);

    \draw[subset] (-3.5,0) -- (3.5,0) node[right] {$\ell$};
    \draw[dashed,line width=0.9pt,draw=blue!65] (P) circle (2.9);
    \draw[dashed,line width=0.9pt,draw=orange!85!black] (A) circle (2.9);
    \draw[dashed,line width=0.9pt,draw=orange!85!black] (B) circle (2.9);
    \draw[very thick] (0,-2.8) -- (0,2.8) node[above] {$m=PQ$};
    \draw (0.28,0) -- (0.28,0.28) -- (0,0.28);

    \fill[geometricSubsetColor] (P) circle (2pt) node[right] {$P$};
    \fill[geometricSubsetColor] (Q) circle (2pt) node[right] {$Q$};
    \fill[geometricSubsetColor] (A) circle (2pt) node[below] {$A$};
    \fill[geometricSubsetColor] (B) circle (2pt) node[below] {$B$};

    \node[blue!65,align=left] at (-3.0,2.65)
        {\small Step 1: circle centred at $P$};
    \node[orange!85!black,align=left] at (3.0,-2.65)
        {\small Step 2: circles centred at $A$ and $B$};
\end{tikzpicture}
\end{center}

\begin{theorembox}
The line $PQ$ is the perpendicular bisector of $AB$ and therefore is
perpendicular to $\ell$ through the prescribed external point $P$.
\end{theorembox}

\subsubsection{Bisecting an Angle}

\begin{derivationbox}[Construction and justification]
Let two rays with common endpoint $O$ form an angle. Draw a circle centred at
$O$ meeting the rays at $A$ and $B$. Then draw equal-radius circles centred at
$A$ and $B$, and let one of their intersections inside the angle be $P$. Since
\[
|OA|=|OB|,
\qquad
|AP|=|BP|,
\qquad
|OP|=|OP|,
\]
the triangles $OAP$ and $OBP$ are congruent.
\end{derivationbox}

\begin{theorembox}
The congruence of $OAP$ and $OBP$ gives
\[
\angle AOP=\angle POB.
\]
Therefore the ray $OP$ bisects the angle.
\end{theorembox}

\subsubsection{Constructing a Parallel Through a Point}

\begin{derivationbox}[Construction]
Let $P$ lie outside a line $\ell$. First construct a line $m$ through $P$
perpendicular to $\ell$. Then construct a line $n$ through $P$ perpendicular
to $m$.

The transversal $m$ forms equal corresponding right angles with $n$ and
$\ell$. By the converse of the corresponding-angles theorem, the two lines are
parallel.
\end{derivationbox}

\begin{theorembox}
The constructed line satisfies
\[
\boxed{n\parallel\ell}.
\]
The fifth postulate, equivalently Playfair's postulate, guarantees that this is
the unique parallel to $\ell$ through $P$.
\end{theorembox}
